% !Mode:: "TeX:UTF-8"
% 文字编码：UTF-8

% \eqnarray如果很长，影响分栏、换行和分页（整块挪动，造成页面空白），
% 可以设置成为自动调整模式
\allowdisplaybreaks[4]

%%%%%%%%%%%%%%%%%%%%%%%%%%%%%%%%%%%%%%%%%%%%%%%%%%%%%%%%%%%
%下面这组命令使浮动对象的缺省值稍微宽松一点，从而防止幅度
%对象占据过多的文本页面，也可以防止在很大空白的浮动页上放置
%很小的图形。
%%%%%%%%%%%%%%%%%%%%%%%%%%%%%%%%%%%%%%%%%%%%%%%%%%%%%%%%%%%
\renewcommand{\textfraction}{0.15}
\renewcommand{\topfraction}{0.85}
\renewcommand{\bottomfraction}{0.65}
\renewcommand{\floatpagefraction}{0.60}

%公式与正文之间的距离
\setlength{\abovedisplayskip}{6pt plus 2pt minus 2pt}
\setlength{\abovedisplayshortskip}{6pt plus 2pt minus 2pt}
\setlength{\belowdisplayskip}{6pt plus 2pt minus 2pt}
\setlength{\belowdisplayshortskip}{6pt plus 2pt minus 2pt}

%%%%%%%%%%%%%%%%%%%%%%%%%%%%%%%%%%%%%%%%%%%%%%%%%%%%%%%%%%%
%下面这组命令可以使公式编号随着每开始新的一节而重新开始。
%%%%%%%%%%%%%%%%%%%%%%%%%%%%%%%%%%%%%%%%%%%%%%%%%%%%%%%%%%%

%\makeatletter      % '@' is now a normail "letter" for TeX
%\@addtoreset{eqation}{section}
%\makeatother       % '@' is restored as a "non-letter" character for TeX

%%%%%%%%%%%%%%%%%%%%%%%%%%%%%%%%%%%%%%%%%%%%%%%%%%%%%%%%%%%
% 重定义字体命令
%%%%%%%%%%%%%%%%%%%%%%%%%%%%%%%%%%%%%%%%%%%%%%%%%%%%%%%%%%%
% 注意win2000,没有 simsun,　最好到网上找一个
% 一些字体是office2000 带的
%%%%%%%%%%%%%%%%%%%%%%%%%%%%%%%%%%%%%%%%%%%%%%%%%%%%%%%%%%%

\setCJKfamilyfont{song}{SimSun}                             %宋体 song
\newcommand{\song}{\CJKfamily{song}}                        % 宋体   (Windows自带simsun.ttf)

\setCJKfamilyfont{fs}{FangSong}                      %仿宋 fs
\newcommand{\fs}{\CJKfamily{fs}}                            %仿宋体 (Windows自带simfang.ttf)

\setCJKfamilyfont{kai}{KaiTi}                        %楷体 kai
\newcommand{\kai}{\CJKfamily{kai}}                        %楷体 (Windows自带simkai.ttf)

\setCJKfamilyfont{yh}{Microsoft YaHei}                    %微软雅黑 yh
\newcommand{\yh}{\CJKfamily{yh}}

\setCJKfamilyfont{hei}{SimHei}                                    %黑体  hei
\newcommand{\hei}{\CJKfamily{hei}}                          % 黑体   (Windows自带simhei.ttf)

%%%%%%%%%%%%%%%%%%%%%%%%%%%%%%%%%%%%%%%%%%%%%%%%%%%%%%%%%%%
% 重定义字号命令
%%%%%%%%%%%%%%%%%%%%%%%%%%%%%%%%%%%%%%%%%%%%%%%%%%%%%%%%%%%


\newcommand{\typethesis}{\zihao{-0}\hei}      % 封面中的“博士研究生学位论文” 黑体 小初
\newcommand{\typetitle}{\fontsize{26pt}{50pt}\hei\bfseries\selectfont}     % 封面中的标题内容 黑体加粗 一号, 1.25倍行距
\newcommand{\typetimu}{\fontsize{22pt}{50pt}\selectfont}     % 封面中的“题目” 宋体 二号, 同一号字体1.25倍行距
\newcommand{\typename}{\fontsize{16pt}{32pt}\fs\selectfont}         % 封面中的姓名内容 仿宋 三号, 1.25倍行距
\newcommand{\typexingming}{\fontsize{15pt}{32pt}\hei\selectfont}    % 封面中的“姓名” 黑体 小三, 同三号字体1.25倍行距
\newcommand{\typedate}{\zihao{3}}         % 封面中的年月 宋体 三号, 1.25倍行距

\newcommand{\typetitleen}{\fontsize{26pt}{48pt}\sffamily\bfseries\itshape\selectfont}     % 英文标题内容 Arial 加粗倾斜 一号, 单倍行距

\newcommand{\typeenglish}{\fontsize{16pt}{22pt}\selectfont}         % 英文封面 三号, 20磅行距
\newcommand{\typeenglishlarge}{\fontsize{16pt}{34pt}\selectfont}         % 英文封面 三号, 20磅行距 加段落间距
\newcommand{\typetitleab}{\fontsize{16pt}{32pt}\sffamily\selectfont}     % 英文摘要页标题内容 Arial 三号, 单倍行距
\newcommand{\typeab}{\fontsize{12pt}{12pt}\sffamily\selectfont}     % 英文摘要页ABSTRACT Arial 小四号, 单倍行距

\newcommand{\typetext}{\fontsize{12pt}{20pt}\selectfont}       % 正文  小四, 20磅
\newcommand{\typechap}{\zihao{3}\linespread{1.64}\hei}         % “摘要”两字 各章标题 黑体 三号, 单倍行距
\newcommand{\typesec}{\zihao{4}\linespread{1.2}\hei}          % 各节标题 黑体 四号, 20磅行距
\newcommand{\typesubsecctex}{\fontsize{13pt}{13pt}\linespread{1.2}\hei} % 二级节标题 黑体 13磅字, 20磅行距
\newcommand{\typesubsubsecctex}{\fontsize{12pt}{12pt}\linespread{1.2}\hei} % 三级节标题 黑体 12磅字, 20磅行距


\newcommand{\typeyemei}{\fontsize{10.5pt}{20pt}\selectfont}      %  页眉 宋体 五号, 20磅行距
\newcommand{\typeyemeien}{\fontsize{10.5pt}{20pt}\selectfont}      %  页眉 英文 五号, 20磅行距
\newcommand{\typeyejiao}{\fontsize{10.5pt}{20pt}\selectfont}          %  页脚 英文 五号, 20磅行距
\newcommand{\typefootnote}{\fontsize{9pt}{15.7pt}\selectfont}          %   脚注 小五号, 单倍行距
\newcommand{\typefootscript}{\fontsize{8pt}{8pt}\selectfont}          %   脚注上标 小五号, 单倍行距
\newcommand{\typetable}{\small\renewcommand\arraystretch{1.5}}          %   表的内容 11磅, 单倍行距
\newcommand{\typebib}{\fontsize{10.5pt}{16pt}\selectfont}      %  参考文献 五号, 16磅行距

%%%%%%%%%%%%%%%%%%%%%%%%%%%%%%%%%%%%%%%%%%%%%%%%%%%%%%%%%%%
% 重定义一些正文相关标题
%%%%%%%%%%%%%%%%%%%%%%%%%%%%%%%%%%%%%%%%%%%%%%%%%%%%%%%%%%%
\CTEXsetup[format={\typechap\centering},nameformat={\typechap},aftername={\hspace*{1em}},titleformat={\typechap},beforeskip={-2pt},afterskip={24pt}]{chapter}
\CTEXsetup[format={\typesec},nameformat={\typesec},aftername={\hspace*{1em}},titleformat={\typesec},beforeskip={24pt},afterskip={6pt}]{section}
\CTEXsetup[format={\typesubsecctex},nameformat={\typesubsecctex},aftername={\hspace*{1em}},titleformat={\typesubsecctex},beforeskip={12pt},afterskip={6pt}]{subsection}
\CTEXsetup[format={\typesubsubsecctex},nameformat={\typesubsubsecctex},aftername={\hspace*{1em}},titleformat={\typesubsubsecctex},beforeskip={12pt},afterskip={6pt}]{subsubsection}



\theoremstyle{plain} \theorembodyfont{\song\rmfamily}
\theoremheaderfont{\hei\rmfamily} \theoremseparator{\ }
\newtheorem{definition}{\hei 定义}[chapter]
\newtheorem{proposition}[definition]{\hei 命题}
\newtheorem{assumption}[definition]{\hei 假设}
\newtheorem{lemma}[definition]{\hei 引理}
\newtheorem{remark}[definition]{\hei 注}
%\newtheorem{theorem}{\hei 定理}[chapter]
\newtheorem{theorem}[definition]{\hei 定理}
\newtheorem{axiom}{\hei 公理}
%\newtheorem{algorithm1}[definition]{\hei 算法}
\newtheorem{corollary}[definition]{\hei 推论}
\newtheorem{exercise}[definition]{}

\theoremheaderfont{\CJKfamily{hei}\rmfamily}\theorembodyfont{\rmfamily}
\theoremstyle{nonumberplain} \theoremseparator{:}
\theoremsymbol{$\blacksquare$}
\newtheorem{proof}{\hei 证明}

\theoremsymbol{$\square$}
\newtheorem{example}{\hei 例}


%%%%%%%%%%%%%%%%%%%%%%%%%%%%%%%%%%%%%%%%%%%%%%%%%%%%%%%%%%
% 脚注字体：宋体小五，单倍行距。悬挂缩进 1.5 字符。标号在正文中是上标，在脚注中为正体
\renewcommand{\thefootnote}{}
\renewcommand{\footnoterule}{\noindent \rule{5cm}{0.6pt} \vspace*{6pt} }
\renewcommand{\footnotesize}{\typefootnote}
\renewcommand{\footnotesep}{11pt}

%使用pifont包里\ding{172}~\ding{181}产生的带圈数字1~10
\newcommand\footnotecircle[1]{\footnote{\hangindent 1.5em \hspace*{-2.2em} \ding{\numexpr171+\value{footnote}}\hspace*{0.4em} #1}\textsuperscript{\typefootscript \ding{\numexpr171+\value{footnote}}}}


%%%%%%%%%%%%%%%%%%%%%%%%%%%%%%%%%%%%%%%%%%%%%%%%%%%%%%%%%%%
% 用于中文段落缩进 和正文版式
%%%%%%%%%%%%%%%%%%%%%%%%%%%%%%%%%%%%%%%%%%%%%%%%%%%%%%%%%%%



%%%%%%%%%%%%%%%%%%%%%%%%%%%%%%%%%%%%%%%%%%%%%%%%%%
%定义段落章节的标题和目录项的格式
%%%%%%%%%%%%%%%%%%%%%%%%%%%%%%%%%%%%%%%%%%%%%%%%%%
\setcounter{secnumdepth}{4}
\setcounter{tocdepth}{2}

% Modified By Lei Wang BaconChina
% THU Version
%\titleformat{\chapter}[hang]
%    {\normalfont\xiaosan\filcenter\hei\sf}
%    {\xiaosan{\chaptertitlename}}
%    {20pt}{\xiaosan}
%%\titlespacing{\chapter}{0pt}{-3ex  plus .1ex minus .2ex}{2.5ex plus .1ex minus .2ex}
%\titlespacing{\chapter}{0pt}{-3ex  plus .1ex minus .2ex}{0.25em}
%
%\titleformat{\section}[hang]{\hei \sf \sihao}
%    {\sihao \thesection}{0.5em}{}{}
%%\titlespacing{\section}{0pt}{1.5ex plus .1ex minus .2ex}{\wordsep}
%\titlespacing{\section}{0pt}{-0.2em}{0.8em}
%
%\titleformat{\subsection}[hang]{\hei \sf \banxiaosi}
%    {\banxiaosi \thesubsection}{0.5em}{}{}
%%    {\banxiaosi \thesubsection}{0pt}{}{}
%%\titlespacing{\subsection}{0pt}{1.5ex plus .1ex minus .2ex}{\wordsep}
%\titlespacing{\subsection}{0pt}{-0.25em}{1em}
%
%\titleformat{\subsubsection}[hang]{\hei \sf}
%    {\thesubsubsection }{0.5em}{}{}
%%\titlespacing{\subsubsection}{0pt}{1.2ex plus .1ex minus .2ex}{\wordsep}
%\titlespacing{\subsubsection}{0pt}{0.25em}{0pt}

%\titleformat{\chapter}[hang]
%    {\normalfont\typechap\filcenter}
%    {\typechap{\chaptertitlename}}
%    {1em}{\typechap}
%\titlespacing{\chapter}{0pt}{}{2.5ex plus .1ex minus .2ex}
%\titlespacing{\chapter}{0pt}{15pt}{27pt}

%\titleformat{\section}[hang]{\hei \sf \xiaoer}
%    {\xiaoer \thesection}{0.5em}{}{}
%%\titlespacing{\section}{0pt}{1.5ex plus .1ex minus .2ex}{\wordsep}
%\titlespacing{\section}{0pt}{-0.2em}{0.8em}

%\titleformat{\subsection}[hang]{\hei \sf \sanhao }
%    {\sanhao \thesubsection}{0.5em}{}{}
%%    {\banxiaosi \thesubsection}{0pt}{}{}
%%\titlespacing{\subsection}{0pt}{1.5ex plus .1ex minus .2ex}{\wordsep}
%\titlespacing{\subsection}{0pt}{-0.25em}{1em}

%\titleformat{\subsubsection}[hang]{\hei \sf \sihao}
%    {\sihao \thesubsubsection }{0.5em}{}{}
%%\titlespacing{\subsubsection}{0pt}{1.2ex plus .1ex minus .2ex}{\wordsep}
%\titlespacing{\subsubsection}{0pt}{0.25em}{0pt}


%\titleformat{\subsubsubsection}[hang]{\hei \sf}
%    {\thesubsubsection }{0.5em}{}{}
%%\titlespacing{\subsubsection}{0pt}{1.2ex plus .1ex minus .2ex}{\wordsep}
%\titlespacing{\subsubsubsection}{0pt}{0.5em}{0pt}

%去掉中间对齐的sectionformat，这样就把节的标题左对齐了。
%\renewcommand \sectionformat{}

% 按清华标准, 缩小目录中各级标题之间的缩进
\dottedcontents{chapter}[4em]{\hei\vspace{6pt}}{4em}{10pt}
\dottedcontents{section}[4.2em]{}{2.2em}{5pt}
\dottedcontents{subsection}[7em]{}{3em}{5pt}
%\dottedcontents{subsubsection}[2.86cm]{}{3.4em}{5pt}
%\dottedcontents{subsubsubsection}[3.66cm]{}{4.2em}{5pt}

%\makeatletter
%\@ifundefined{chapter}{}{
%  \newcommand\CJKprechaptername{第}
%  \newcommand\CJKchaptername{章}
%  \newcommand\CJKthechapter{\zhnumber{\@arabic\c@chapter}}}

%\renewcommand\appendixname{附录\@Alph\c@chapter}

%\@ifundefined{mainmatter}
%  {\renewcommand\abstractname{摘要}}{}

%\renewcommand\chaptername{\CJKprechaptername\CJKthechapter\CJKchaptername}
%\makeatother


%%%%%%%%%%%%%%%%%%%%%%%%%%%%%%%%%%%%%%%%%%%%%%%%%%%%%%%
% 定义页眉和页脚 使用fancyhdr 宏包
%%%%%%%%%%%%%%%%%%%%%%%%%%%%%%%%%%%%%%%%%%%%%%%%%%%%%%%%

\newcommand{\makeheadrule}{%
   % \makebox[0pt][l]{\rule[.7\baselineskip]{\headwidth}{0.75pt}}}%\vskip-.8\baselineskip}
   \hrule width\headwidth height0.75pt}



\pagestyle{fancyplain}

%去掉章节标题中的数字
%\renewcommand{\chaptermark}[1]{\markboth{\chaptername\hspace{1em}#1}{}}

\fancyhf{}


%%% Clear Header %%%%%%%%%%%%%%%%%%%%%%%%%%%%%%%%%%%%%%%%%%%%%%%%%%%
% Clear Header Style on the Last Empty Odd pages
\makeatletter
\def\cleardoublepage{\clearpage\if@twoside \ifodd\c@page\else%
    \hbox{}%
    \thispagestyle{empty}%              % Empty header styles
    \newpage%
    \if@twocolumn\hbox{}\newpage\fi\fi\fi}

%%%%%%%%%%%%%%%%%%%%%%%%%%%%%%%%%%%%%%%%%%%%%%%%%%%%%%%%
% 设置行距和段落间垂直距离
%%%%%%%%%%%%%%%%%%%%%%%%%%%%%%%%%%%%%%%%%%%%%%%%%%%%%%%%

% 段落之间的竖直距离
\setlength{\parskip}{0pt}

% 定义行距
\renewcommand{\baselinestretch}{1}

%%%%%%%%%%%%%%%%%%%%%%%%%%%%%%%%%%%%%%%%%%%%%%%%%%%%%%%%
% 调整列表环境的垂直间距
%%%%%%%%%%%%%%%%%%%%%%%%%%%%%%%%%%%%%%%%%%%%%%%%%%%%%%%%
\let\orig@Itemize =\itemize
\let\orig@Enumerate =\enumerate
\let\orig@Description =\description

\def\Myspacing{\itemsep=2ex \topsep=-4ex \partopsep=-2ex \parskip=-1ex \parsep=2ex}

\def\newitemsep{
\renewenvironment{itemize}{\orig@Itemize\Myspacing}{\endlist}
\renewenvironment{enumerate}{\orig@Enumerate\Myspacing}{\endlist}
\renewenvironment{description}{\orig@Description\Myspacing}{\endlist}
}

\def\olditemsep{
\renewenvironment{itemize}{\orig@Itemize}{\endlist}
\renewenvironment{enumerate}{\orig@Enumerate}{\endlist}
\renewenvironment{description}{\orig@Description}{\endlist}
}

\newitemsep

%%%%%%%%%%%%%%%%%%%%%%%%%%%%%%%%%%%%%%%%%%%%%%%%%%%%%%%
% 修改引用的格式，
%%%%%%%%%%%%%%%%%%%%%%%%%%%%%%%%%%%%%%%%%%%%%%%%%%%%%%%

%第一行在引用处数字两边加方框
%第二行去除参考文献里数字两边的方框
%\makeatletter
%\def\@cite#1{\mbox{$\m@th^{\hbox{\@ove@rcfont[#1]}}$}}
%\renewcommand\@biblabel[1]{#1}
%\makeatother

\renewcommand{\@openbib@code}{\addtolength{\itemsep}{-7pt}}

% 增加 \ucite 命令使显示的引用为上标形式
\newcommand{\ucite}[1]{\textsuperscript{\mbox{\scriptsize \cite{#1}}}}%张宏剑修改20130327

%%%%%%%%%%%%%%%%%%%%%%%%%%%%%%%%%%%%%%%%%%%%%%%%%%%%%%%%%%%
%
% 定制浮动图形和表格标题样式
%
%%%%%%%%%%%%%%%%%%%%%%%%%%%%%%%%%%%%%%%%%%%%%%%%%%%%%%%%%%%

%\renewcommand{\captionfont}{\CJKfamily{song}\rmfamily}
%\renewcommand{\captionlabelfont}{\CJKfamily{song}\rmfamily}
%三线表设置
\setlength{\heavyrulewidth}{2.25pt}
\setlength{\lightrulewidth}{0.75pt}
\setlength{\abovetopsep}{0pt}
\setlength{\belowbottomsep}{-6pt}

% 按清华标准, 去掉图表号后面的:
%\renewcommand{\captionlabeldelim}{\hspace{1em}}
%\DeclareCaptionLabelSeparator{colon}{\hspace{1em}}
%\renewcommand{\arraystretch}{1}


% \subsubsection{浮动对象以及表格}
% \label{sec:float}
% 设置浮动对象和文字之间的距离
% \changes{v2.6}{2006/06/09}{增加 \cs{floatsep}，\cs{@fptop}，\cs{@fpsep} 和 \cs{@fpbot}。}
%    \begin{macrocode}
%<*cls>
%\makeatletter
%\setlength{\floatsep}{12bp \@plus4pt \@minus1pt}
%\setlength{\intextsep}{2pt \@plus2pt \@minus2pt}
%\setlength{\textfloatsep}{12bp \@plus4pt \@minus2pt}
%\setlength{\@fptop}{0bp \@plus1.0fil}
%\setlength{\@fpsep}{12bp \@plus2.0fil}
%\setlength{\@fpbot}{0bp \@plus1.0fil}
%\makeatother
%    \end{macrocode}

% 定制浮动图形和表格标题样式（参考ThuThesis：清华大学学位论文模板）
% \begin{itemize}
%   \item 图表标题字体为 11pt， 这里写作大五号
%   \item 去掉图表号后面的冒号。图序与图名文字之间空一个汉字符宽度。
%   \item 图：caption 在下，段前空 6 磅，段后空 12 磅
%   \item 表：caption 在上，段前空 12 磅，段后空 6 磅
% \end{itemize}

\captionsetup[table]{labelsep=quad,format=hang,justification=centerlast,font={small,stretch=1.16},position=top,belowskip={15pt-\intextsep},aboveskip=6pt}

\captionsetup[figure]{labelsep=quad,format=hang,justification=centerlast,font={small,stretch=1.16},position=bottom,belowskip={10pt-\intextsep},aboveskip=6pt}

\captionsetup[subfigure]{labelsep=quad,format=hang,justification=centerlast,font={small,stretch=1.16},position=bottom,belowskip=0pt,aboveskip=6pt}


%%%%%%%%%%%%%%%%%%%%%%%%%%%%%%%%%%%%%%%%%%%%%%%%%%%%%%%
% 定义题头格言的格式
%%%%%%%%%%%%%%%%%%%%%%%%%%%%%%%%%%%%%%%%%%%%%%%%%%%%%%%

%
% 用法 \begin{Aphorism}{author}
%         aphorism
%      \end{Aphorism}

\newsavebox{\AphorismAuthor}
\newenvironment{Aphorism}[1]
{\vspace{0.5cm}\begin{sloppypar} \slshape
\sbox{\AphorismAuthor}{#1}
\begin{quote}\small\itshape }
{\\ \hspace*{\fill}------\hspace{0.2cm} \usebox{\AphorismAuthor}
\end{quote}
\end{sloppypar}\vspace{0.5cm}}

%自定义一个空命令，用于注释掉文本中不需要的部分。
\newcommand{\comment}[1]{}

% This is the flag for longer version
\newcommand{\longer}[2]{#1}

\newcommand{\ds}{\displaystyle}

% define graph scale
\def\gs{1.0}

%%%%%%%%%%%%%%%%%%%%%%%%%%%%%%%%%%%%%%%%%%%%%%%%%%%%%%%%%%%%%%%%%%%%%%
% 自定义项目列表标签及格式 \begin{denselist} 列表项 \end{denselist}
%%%%%%%%%%%%%%%%%%%%%%%%%%%%%%%%%%%%%%%%%%%%%%%%%%%%%%%%%%%%%%%%%%%%%%
\newcounter{newlist} %自定义新计数器
\newenvironment{denselist}[1][temp]{%%%%% 定义新环境：可改变的列表题目
\begin{list}{\textbf{#1} (\arabic{newlist})} %%标签格式
    {
    \usecounter{newlist}
     \setlength{\labelwidth}{22pt} %标签盒子宽度
     \setlength{\labelsep}{0cm} %标签与列表文本距离
     \setlength{\leftmargin}{0cm} %左右边界
     \setlength{\rightmargin}{0cm}
     \setlength{\parsep}{0ex} %段落间距
     \setlength{\itemsep}{0ex} %标签间距
     \setlength{\itemindent}{44pt} %标签缩进量
     \setlength{\listparindent}{44pt} %段落缩进量
    }}
{\end{list}}%%%%%

%\def\defaultfont{\renewcommand{\baselinestretch}{1.5}
%\fontsize{12pt}{13pt}\selectfont}

%%%%%%%%%%%%%%%%%%%%%%%%%%%%%%%%%%%%%%%%%%%%%%%%%%%%%%%%%%%%%%%%%%%%%%
% 封面、摘要、版权、致谢格式定义 --- by Feng Hua & Lei Wang
%%%%%%%%%%%%%%%%%%%%%%%%%%%%%%%%%%%%%%%%%%%%%%%%%%%%%%%%%%%%%%%%%%%%%%

\def\ctitlefirstline#1{\def\@ctitlefirstline{#1}}\def\@ctitlefirstline{}
\def\ctitlesecondline#1{\def\@ctitlesecondline{#1}}\def\@ctitlesecondline{}
\def\cdegree#1{\def\@cdegree{#1}}\def\@cdegree{}
\def\caffil#1{\def\@caffil{#1}}\def\@caffil{}
\def\csubject#1{\def\@csubject{#1}}\def\@csubject{}
\def\cauthor#1{\def\@cauthor{#1}}\def\@cauthor{}
\def\cnumber#1{\def\@cnumber{#1}}\def\@cnumber{}
\def\csubsub#1{\def\@csubsub{#1}}\def\@csubsub{}
\def\csupervisor#1{\def\@csupervisor{#1}}\def\@csupervisor{}
\def\cassosupervisor#1{\def\@cassosupervisor{~ & 副指导教师 & ：& #1\\}}\def\@cassosupervisor{}
\def\ccosupervisor#1{\def\@ccosupervisor{~ & 联\hfill合\hfill导\hfill 师 & ：& #1\\}}\def\@ccosupervisor{}
\def\cdate#1{\def\@cdate{#1}}\def\@cdate{}
\long\def\cabstract#1{\long\def\@cabstract{#1}}\long\def\@cabstract{}
\def\ckeywords#1{\def\@ckeywords{#1}}\def\@ckeywords{}

\def\etitle#1{\def\@etitle{#1}}\def\@etitle{}
\def\edegree#1{\def\@edegree{#1}}\def\@edegree{}
\def\eaffil#1{\def\@eaffil{#1}}\def\@eaffil{}
\def\esubject#1{\def\@esubject{#1}}\def\@esubject{}
\def\edepartment#1{\def\@edepartment{#1}}\def\@edepartment{}
\def\eauthor#1{\def\@eauthor{#1}}\def\@eauthor{}
\def\esupervisor#1{\def\@esupervisor{#1}}\def\@esupervisor{}
\def\eassosupervisor#1{\def\@eassosupervisor{Associate Supervisor: & #1\\}}\def\@eassosupervisor{}
\def\ecosupervisor#1{\def\@ecosupervisor{~ & #1\\}}\def\@ecosupervisor{}
\def\edate#1{\def\@edate{#1}}\def\@edate{}
\long\def\eabstract#1{\long\def\@eabstract{#1}}\long\def\@eabstract{}
\def\ekeywords#1{\def\@ekeywords{#1}}\def\@ekeywords{}

\def\makecover{
    \begin{titlepage}
    %中文封面%%%%%%%%%%%%%%%%%%%%%%%%%%%%%%%%%%%%%%%%%%%%%%%%%%%
    \newpage
    \thispagestyle{empty}
    \begin{center} {
    \vspace*{3mm}
%    \centerline{\epsfig{file=newlogo.eps,width=0.7\columnwidth,clip=}}%张宏剑
%    \epsfig{file=logo.eps,width=0.53\columnwidth,clip=} \\
    \includegraphics[width=8.5cm]{logo.png} \\
    \vspace*{4mm}
    {\typethesis \@cdegree 研究生学位论文} \\
    \vspace*{8mm}
    {\typetimu \hspace*{-1cm}
    \begin{tabular}{lc}
    \raisebox{-2mm}{题目：} & {\typetitle \@ctitlefirstline} \\ \cline{2-2}
          & {\typetitle \@ctitlesecondline} \\ \cline{2-2}
    　 　 & {\hspace*{9cm}}　
    \end{tabular} } \\
    \vspace*{2.2cm}
    {\typexingming \hspace*{-2cm}
    \begin{tabular}{lc}
    姓\hphantom{姓名}名： \hspace*{1mm}& {\typename \@cauthor}\\ \cline{2-2}
    学\hphantom{姓名}号： & {\typename \@cnumber}\\ \cline{2-2}
    院\hphantom{姓名}系： & {\typename \@caffil}\\ \cline{2-2}
    专\hphantom{姓名}业： & {\typename \@csubject}\\ \cline{2-2}
    研究方向： & {\typename \@csubsub}\\ \cline{2-2}
    导师姓名： & {\typename \@csupervisor} \\ \cline{2-2}
    　 　 & {\hspace*{6.5cm}}　
    \end{tabular} } \\
    \vspace*{1.2cm}
    \typedate \@cdate
    }
    \end{center}

    % 封二 空白页
    \cleardoublepage
     \newpage
    \thispagestyle{empty}
    \begin{center}
    \vspace*{-1.2cm}
    {\typetitleen
    \begin{center} { \@etitle} \end{center}}
    \vspace*{1.8cm}
    {\typeenglish
    Dissertation Submitted to \\
    \textbf{Peking University} \\
    \vspace*{0.6cm}
    \textbf{In partial fulfillment of the requirements for the degree of} \\
    \@edegree \\
    \vspace*{2cm}
    By \\
    \textbf{\@eauthor , Master Candidate} \\
    \textbf{(\@esubject)} \\}
    \vspace*{0.8cm}
    {\typeenglishlarge
     Dissertation Supervisor: \textbf{\@esupervisor} \\ \quad \\
    Department of \@edepartment \\
    College of Electronic Engineering and Computer Science \\
    Peking University \\
    \@edate, Beijing \\}
    \end{center}

    \cleardoublepage
     \newpage
    \thispagestyle{empty}
    \end{titlepage}
}

\def\makeabstract{
   \begin{titlepage}%%%张宏剑
   \newpage%张宏剑
    \thispagestyle{plain}%张宏剑
    %Abstract and keywords%%%%%%%%%%%%%%%%%%%%%%%%%%%%%%%%%%%%%%%%%%%%%%%%%%%%%
    %\newpage%张宏剑施加试试不出现yemei
    %\thispagestyle{empty}%张宏剑施加试试不出现yemei
%    \defaultfont

    \chapter*{摘要 \markboth{摘要}{摘要}}
    %\addcontentsline{toc}{chapter}{\hei {摘~~~~要}}%陈立洋注释掉：2014年5月更改，目录从第一章开始，不包含摘要
    %张宏剑去掉
    \setcounter{page}{1}
    {\typetext \@cabstract }

    \cleardoublepage
    \newpage

    \thispagestyle{plain}%页眉上不出现文字
%    \chapter*{ABSTRACT}
    %\addcontentsline{toc}{chapter}{\hei Abstract} %陈立洋注释掉：2014年5月学院要求，目录从第一章开始，不包含摘要
    \markboth{\typeyemeien ABSTRACT}{\typeyemeien ABSTRACT}%张宏剑去掉
    \vspace*{-5mm}
    \begin{center}
    {\typetitleab
    \begin{center} {\@etitle}\end{center}}
    \vspace*{8mm}
    {\typetext
    \@eauthor\ (\@esubject) \\
    Directed by \@esupervisor \\[3mm] }
    {\typeab \textbf{ABSTRACT}}
    \end{center}
    \vspace*{-2mm}
    \typetext
    \@eabstract
%    \noindent {\textbf{Key Words:}} \quad \@ekeywords
    \cleardoublepage

     \end{titlepage}%%张宏剑
}
\makeatother



\raggedbottom

%\oddsidemargin 0.9em

\setlength{\multlinegap}{0mm}
